\section{实验对照与结论安排}

各项实验围绕总体有效性、偏离状态适应和领域适用性展开。主要训练对照使用同一起点的普通组A与恢复组C，两组匹配问题和监督token；输入历史长度及计算开销单独记录。雷达域内实验采用同样的对照原则，$I_{\mathrm{domain}}$用于展示适配起点，$A_{\mathrm{domain}}$作为恢复训练增量的主要基线。

\begin{table}[htbp]
\centering\footnotesize
\caption{主要实验的研究问题、对照与状态}
\label{tab:experiment-matrix}
\begin{tabular}{p{2.1cm}p{3.9cm}p{4.4cm}p{2cm}}
\toprule
研究问题 & 固定材料与对照 & 主要指标与结论对象 & 状态 \\
\midrule
总体问答效果 & KQA Pro项目留出2000题；两次配对续训A/C，P作参照 & 准确率、完成率与成本；A/C平均78.75\%/83.50\% & 已完成 \\
偏离历史适应 & 固定2400题池中的合格状态；同一正常／偏离历史下A/C & 偏离状态的C减A增量，以及相对正常状态的增量变化 & 待推理 \\
条件证据组织 & 八来源96题；相同语义记录的平铺／绑定 & 正文正确73/80；绑定错误17/11；引用支持另列 & 已完成 \\
领域恢复作用 & 同一$I_{\mathrm{domain}}$、工具与目录；普通／恢复后缀续训 & 新来源基础及组合题分别报告完整答案、来源分组、完成率和成本 & 任务已准备；训练待开展 \\
\bottomrule
\end{tabular}
\end{table}

公开实验给出恢复监督的总体性能证据。同状态评价将两组模型置于一致的实际历史中：偏离条件下的正向差异说明该类状态中的求解收益，偏离条件相对正常条件更大的增量进一步支持对偏离历史的适应。领域配对实验则依据完整参数答案评价同一方法在雷达任务中的作用。条件绑定实验独立刻画知识组织对领域回答的影响。已完成的强化学习扩展和反馈文字诊断作为方法分析材料纳入第五章。
